\chapter*{Synthèse des formules essentielles}
\markboth{Synthèse des formules essentielles}{Synthèse des formules essentielles}
\addcontentsline{toc}{chapter}{Synthèse des formules essentielles}

\section*{Algèbre et fonctions}
\begin{align*}
(a+b)^2&=a^2+2ab+b^2,
&a^2-b^2&=(a-b)(a+b),\\
x_{1,2}&=\frac{-b\pm\sqrt{b^2-4ac}}{2a},
&\log_bx&=\frac{\ln x}{\ln b}.
\end{align*}

Les racines quadratiques supposent \(a\ne0\) et \(b^2-4ac\ge0\) pour rester réelles. Pour le logarithme, \(b>0\), \(b\ne1\) et \(x>0\).

\section*{Dénombrement et probabilités}
\begin{align*}
P_n&=n!,&
A_n^k&=\frac{n!}{(n-k)!},\\
\binom nk&=\frac{n!}{k!(n-k)!},&
P(A^c)&=1-P(A),\\
P(A\mid B)&=\frac{P(A\cap B)}{P(B)},&
P(A\cup B)&=P(A)+P(B)-P(A\cap B).
\end{align*}

Pour les arrangements et combinaisons, \(n,k\) sont entiers et \(0\le k\le n\). La probabilité conditionnelle suppose \(P(B)>0\).

\section*{Vecteurs et matrices}
\begin{align*}
\norm{\mathbf v}&=\sqrt{v_1^2+\cdots+v_n^2},
&\mathbf u\cdot\mathbf v&=\norm{\mathbf u}\norm{\mathbf v}\cos\theta,\\
\det\begin{pmatrix}a&b\\c&d\end{pmatrix}&=ad-bc,
&\begin{pmatrix}a&b\\c&d\end{pmatrix}^{-1}
&=\frac1{ad-bc}\begin{pmatrix}d&-b\\-c&a\end{pmatrix}.
\end{align*}

L'angle \(\theta\) exige deux vecteurs non nuls; l'inverse de la matrice exige \(ad-bc\ne0\).

\section*{Trigonométrie}
\begin{align*}
\sin^2\theta+\cos^2\theta&=1,
&\tan\theta&=\frac{\sin\theta}{\cos\theta},\\
\frac{a}{\sin A}&=\frac{b}{\sin B}=\frac{c}{\sin C},
&c^2&=a^2+b^2-2ab\cos C,\\
T&=\frac{2\pi}{\abs B},
&|\omega|&=\frac{2\pi}{T}.
\end{align*}

La tangente exige \(\cos\theta\ne0\). Les lois du triangle s'appliquent à un triangle non dégénéré. La période de \(A\sin(B(x-C))+D\) suppose \(A\ne0\) et \(B\ne0\); pour le mouvement circulaire, \(T>0\).

\chapter*{Conclusion}
\markboth{Conclusion}{Conclusion}
\addcontentsline{toc}{chapter}{Conclusion}
Les chapitres de MAT0339 sont reliés par une même idée : choisir une représentation adaptée. Les nombres donnent les objets de base; l'algèbre transforme les expressions; les fonctions décrivent la dépendance; les probabilités modélisent l'incertitude; les vecteurs et matrices organisent l'espace et les systèmes; la trigonométrie décrit les angles, les triangles et les phénomènes périodiques.

Le manuel a volontairement approfondi ces contenus à l'intérieur de leurs chapitres. Les exemples détaillés, les démonstrations, les erreurs fréquentes et les méthodes de vérification ne forment pas un programme supplémentaire : ils rendent le programme principal plus compréhensible et plus autonome.


